\section{System Architecture}
\label{sec:architecture}

\subsection{Design Overview}

Our distributed MPI WiFi architecture follows a centralized-channel, distributed-devices pattern. Figure~\ref{fig:architecture-overview} illustrates the high-level design:

\begin{figure}[htbp]
\centering
\begin{verbatim}
┌──────────────────────────────────────────────────────────────┐
│                     RANK 0: Channel Processor                 │
│  ┌────────────────────────────────────────────────────────┐  │
│  │         WifiChannelMpiProcessor                        │  │
│  │  - Maintains global device registry                    │  │
│  │  - Receives TX_REQUEST messages                        │  │
│  │  - Calculates propagation (loss + delay)               │  │
│  │  - Sends RX_NOTIFICATION to receivers                  │  │
│  └────────────────────────────────────────────────────────┘  │
│         ▲                │                │          ▲        │
│         │ TX_REQUEST     │ RX_NOTIFY      │          │        │
└─────────┼────────────────┼────────────────┼──────────┼────────┘
          │                │                │          │
    ┌─────┴─────┐    ┌─────▼──────┐   ┌────▼─────┐  │
    │           │    │            │   │          │  │
┌───▼───────────▼────▼────────────▼───▼──────────▼──▼─────────┐
│              RANKS 1-N: Device Processors                     │
│  ┌──────────────┐  ┌──────────────┐  ┌──────────────┐       │
│  │  Node 1      │  │  Node 2      │  │  Node n      │       │
│  │ ┌──────────┐ │  │ ┌──────────┐ │  │ ┌──────────┐ │       │
│  │ │WiFi App  │ │  │ │WiFi App  │ │  │ │WiFi App  │ │       │
│  │ └────┬─────┘ │  │ └────┬─────┘ │  │ └────┬─────┘ │       │
│  │ ┌────▼─────┐ │  │ ┌────▼─────┐ │  │ ┌────▼─────┐ │       │
│  │ │WiFi MAC  │ │  │ │WiFi MAC  │ │  │ │WiFi MAC  │ │       │
│  │ └────┬─────┘ │  │ └────┬─────┘ │  │ └────┬─────┘ │       │
│  │ ┌────▼─────┐ │  │ ┌────▼─────┐ │  │ ┌────▼─────┐ │       │
│  │ │WiFi PHY  │ │  │ │WiFi PHY  │ │  │ │WiFi PHY  │ │       │
│  │ └────┬─────┘ │  │ └────┬─────┘ │  │ └────┬─────┘ │       │
│  │ ┌────▼───────────────────────────────────▼─────┐ │       │
│  │ │  RemoteYansWifiChannelStub (MPI Proxy)       │ │       │
│  │ │  - Intercepts Send() calls                   │ │       │
│  │ │  - Serializes to MPI messages                │ │       │
│  │ │  - Receives RX notifications                 │ │       │
│  │ └──────────────────────────────────────────────┘ │       │
│  └──────────────┘  └──────────────┘  └──────────────┘       │
└───────────────────────────────────────────────────────────────┘
\end{verbatim}
\caption{Distributed MPI WiFi architecture overview}
\label{fig:architecture-overview}
\end{figure}

\textbf{Design Rationale}: Rather than attempting to distribute the broadcast channel across machines (which would require complex consistency protocols), we embrace WiFi's inherent centralization. Rank 0 maintains the authoritative channel state, while device simulation—which includes protocol stack processing, application logic, and local state—is distributed across ranks 1 through N. This design:

\begin{itemize}
    \item Maintains channel correctness (single source of truth)
    \item Enables device parallelism (independent protocol processing)
    \item Simplifies synchronization (centralized decisions)
    \item Scales well when computation is device-dominated rather than channel-dominated
\end{itemize}

\subsection{Core Components}

\subsubsection{RemoteYansWifiChannelStub}

The \texttt{RemoteYansWifiChannelStub} class acts as a local proxy for the remote channel processor. It inherits from \texttt{YansWifiChannel} to maintain API compatibility, but overrides key methods to use MPI communication:

\begin{lstlisting}[style=cpp, caption={RemoteYansWifiChannelStub class declaration}]
class RemoteYansWifiChannelStub : public YansWifiChannel
{
public:
    RemoteYansWifiChannelStub();
    virtual ~RemoteYansWifiChannelStub();
    
    // Configuration
    void SetRemoteChannelRank(uint32_t rank);
    void SetLocalDeviceRank(uint32_t rank);
    
    // Override YansWifiChannel methods
    void Send(Ptr<YansWifiPhy> sender,
              Ptr<const WifiPpdu> ppdu,
              double txPowerDbm) override;
    
    void Add(Ptr<YansWifiPhy> phy) override;
    
private:
    void SendDeviceRegistration(Ptr<YansWifiPhy> phy);
    void SendTransmissionRequest(Ptr<YansWifiPhy> sender,
                                 Ptr<const WifiPpdu> ppdu,
                                 double txPowerDbm);
    void ReceiveNotifications(); // MPI receive handler
    
    uint32_t m_channelRank;  // Rank 0
    uint32_t m_localRank;     // This rank's ID
    std::map<uint32_t, Ptr<YansWifiPhy>> m_phyMap;
};
\end{lstlisting}

\paragraph{Device Registration}
When a WiFi PHY is added via \texttt{Add()}, the stub sends a \texttt{DEVICE\_REGISTRATION} message to rank 0:

\begin{lstlisting}[style=cpp, caption={Device registration flow}]
void RemoteYansWifiChannelStub::Add(Ptr<YansWifiPhy> phy)
{
    uint32_t deviceId = m_phyMap.size();
    m_phyMap[deviceId] = phy;
    
    // Get device position
    Ptr<MobilityModel> mobility = phy->GetMobility();
    Vector position = mobility->GetPosition();
    
    // Prepare MPI message
    struct DeviceRegMsg {
        uint32_t srcRank;
        uint32_t deviceId;
        double posX, posY, posZ;
    } msg;
    
    msg.srcRank = m_localRank;
    msg.deviceId = deviceId;
    msg.posX = position.x;
    msg.posY = position.y;
    msg.posZ = position.z;
    
    // Send to channel processor
    MPI_Send(&msg, sizeof(msg), MPI_BYTE,
             m_channelRank, MSG_DEVICE_REG,
             MPI_COMM_WORLD);
             
    NS_LOG_INFO("[RemoteYansWifiChannelStub] Device " 
                << deviceId << " registered via MPI");
}
\end{lstlisting}

\paragraph{Transmission Request}
When a device transmits, \texttt{Send()} serializes the packet and metadata to MPI:

\begin{lstlisting}[style=cpp, caption={Transmission request serialization}]
void RemoteYansWifiChannelStub::Send(
    Ptr<YansWifiPhy> sender,
    Ptr<const WifiPpdu> ppdu,
    double txPowerDbm)
{
    // Find sender device ID
    uint32_t senderId = FindDeviceId(sender);
    
    // Get current position (for mobile devices)
    Ptr<MobilityModel> mobility = sender->GetMobility();
    Vector position = mobility->GetPosition();
    
    // Serialize packet
    Ptr<Packet> packet = ppdu->GetPacket()->Copy();
    uint32_t packetSize = packet->GetSize();
    uint8_t *buffer = new uint8_t[packetSize];
    packet->CopyData(buffer, packetSize);
    
    // Prepare TX request message
    struct TxRequestMsg {
        uint32_t srcRank;
        uint32_t deviceId;
        double txPowerDbm;
        double posX, posY, posZ;
        uint64_t timestamp; // Simulation time
        uint32_t packetSize;
        // Followed by packet data
    } msg;
    
    msg.srcRank = m_localRank;
    msg.deviceId = senderId;
    msg.txPowerDbm = txPowerDbm;
    msg.posX = position.x;
    msg.posY = position.y;
    msg.posZ = position.z;
    msg.timestamp = Simulator::Now().GetNanoSeconds();
    msg.packetSize = packetSize;
    
    // Send header + payload
    MPI_Send(&msg, sizeof(msg), MPI_BYTE,
             m_channelRank, MSG_TX_REQUEST,
             MPI_COMM_WORLD);
    MPI_Send(buffer, packetSize, MPI_BYTE,
             m_channelRank, MSG_TX_PAYLOAD,
             MPI_COMM_WORLD);
    
    delete[] buffer;
    
    NS_LOG_INFO("[RemoteYansWifiChannelStub] TX request sent: "
                << "Device " << senderId 
                << ", Power " << txPowerDbm << " dBm");
}
\end{lstlisting}

\subsubsection{WifiChannelMpiProcessor}

The \texttt{WifiChannelMpiProcessor} runs on rank 0 and handles all channel computations:

\begin{lstlisting}[style=cpp, caption={WifiChannelMpiProcessor class declaration}]
class WifiChannelMpiProcessor : public Object
{
public:
    WifiChannelMpiProcessor();
    virtual ~WifiChannelMpiProcessor();
    
    void SetPropagationLossModel(
        Ptr<PropagationLossModel> loss);
    void SetPropagationDelayModel(
        Ptr<PropagationDelayModel> delay);
    
    void Run(); // Main processing loop
    
private:
    void HandleDeviceRegistration();
    void HandleTransmissionRequest();
    void ProcessTransmission(TxRequestMsg& tx);
    void SendRxNotification(uint32_t destRank,
                           uint32_t destDevice,
                           const Packet& pkt,
                           double rxPowerDbm,
                           Time delay);
    
    struct DeviceInfo {
        uint32_t rank;
        uint32_t deviceId;
        Vector position;
    };
    
    std::map<std::pair<uint32_t,uint32_t>, DeviceInfo> 
        m_devices; // Key: (rank, deviceId)
    
    Ptr<PropagationLossModel> m_lossModel;
    Ptr<PropagationDelayModel> m_delayModel;
};
\end{lstlisting}

\paragraph{Processing Loop}
The channel processor runs a message-processing loop:

\begin{lstlisting}[style=cpp, caption={Channel processor main loop}]
void WifiChannelMpiProcessor::Run()
{
    NS_LOG_INFO("WifiChannelMpiProcessor started on rank 0");
    
    while (Simulator::Now() < Simulator::GetMaximumSimulationTime())
    {
        // Check for incoming MPI messages (non-blocking)
        MPI_Status status;
        int flag;
        MPI_Iprobe(MPI_ANY_SOURCE, MPI_ANY_TAG,
                   MPI_COMM_WORLD, &flag, &status);
        
        if (flag) {
            switch (status.MPI_TAG) {
                case MSG_DEVICE_REG:
                    HandleDeviceRegistration();
                    break;
                    
                case MSG_TX_REQUEST:
                    HandleTransmissionRequest();
                    break;
            }
        }
        
        // Process next simulation event
        Simulator::RunOneEvent();
    }
}
\end{lstlisting}

\paragraph{Transmission Processing}
For each transmission request, the processor calculates propagation to all registered devices:

\begin{lstlisting}[style=cpp, caption={Transmission processing and propagation}]
void WifiChannelMpiProcessor::ProcessTransmission(
    TxRequestMsg& tx)
{
    // Sender info
    Vector txPos(tx.posX, tx.posY, tx.posZ);
    auto senderKey = std::make_pair(tx.srcRank, tx.deviceId);
    
    NS_LOG_INFO("Processing TX from rank " << tx.srcRank
                << ", device " << tx.deviceId
                << " at power " << tx.txPowerDbm << " dBm");
    
    // Calculate propagation to all receivers
    for (auto& [key, rxDevice] : m_devices)
    {
        // Skip sender
        if (key == senderKey) continue;
        
        Vector rxPos = rxDevice.position;
        
        // Calculate path loss
        double rxPowerDbm = m_lossModel->CalcRxPower(
            tx.txPowerDbm, txPos, rxPos);
        
        // Calculate propagation delay
        Time delay = m_delayModel->GetDelay(txPos, rxPos);
        
        // Send RX notification to receiver's rank
        SendRxNotification(rxDevice.rank,
                          rxDevice.deviceId,
                          tx.packet,
                          rxPowerDbm,
                          delay);
        
        NS_LOG_DEBUG("  -> Receiver rank " << rxDevice.rank
                     << ", device " << rxDevice.deviceId
                     << ": RX power " << rxPowerDbm << " dBm, "
                     << "delay " << delay.GetMicroSeconds() << " μs");
    }
}
\end{lstlisting}

\subsection{MPI Message Protocol}

Our protocol uses three primary message types, detailed in Table~\ref{tab:mpi-messages}:

\begin{table}[htbp]
\centering
\caption{MPI message types and payloads}
\label{tab:mpi-messages}
\begin{tabular}{@{}llp{6cm}@{}}
\toprule
\textbf{Message Type} & \textbf{Direction} & \textbf{Payload} \\
\midrule
\texttt{DEVICE\_REG} & Device→Channel & Source rank, device ID, position (x,y,z) \\
\texttt{TX\_REQUEST} & Device→Channel & Source rank, device ID, TX power (dBm), position, timestamp, packet size, packet data \\
\texttt{RX\_NOTIFICATION} & Channel→Device & Destination rank, device ID, RX power (dBm), propagation delay, packet size, packet data \\
\texttt{MOBILITY\_UPDATE} & Device→Channel & Source rank, device ID, new position (x,y,z), velocity \\
\bottomrule
\end{tabular}
\end{table}

\subsection{Communication Flow}

Figure~\ref{fig:sequence-diagram} shows the complete sequence for a packet transmission:

\begin{figure}[htbp]
\centering
\begin{verbatim}
Device (Rank 1)         Stub (Rank 1)         Processor (Rank 0)      Stub (Rank 2)      Device (Rank 2)
     │                       │                        │                      │                  │
     │ Send()                │                        │                      │                  │
     ├──────────────────────>│                        │                      │                  │
     │                       │ TX_REQUEST             │                      │                  │
     │                       ├───────────────────────>│                      │                  │
     │                       │  (packet + metadata)   │                      │                  │
     │                       │                        │                      │                  │
     │                       │                   Calc propagation            │                  │
     │                       │                   to all receivers            │                  │
     │                       │                        │                      │                  │
     │                       │                        │ RX_NOTIFICATION      │                  │
     │                       │                        ├─────────────────────>│                  │
     │                       │                        │  (packet + RX power) │                  │
     │                       │                        │                      │                  │
     │                       │                        │                      │ StartReceive()   │
     │                       │                        │                      ├─────────────────>│
     │                       │                        │                      │                  │
     │                       │                        │                      │                  │ MAC processing
     │                       │                        │                      │                  │ Frame reception
     │                       │                        │                      │                  │ Deliver to app
\end{verbatim}
\caption{Sequence diagram for packet transmission}
\label{fig:sequence-diagram}
\end{figure}

\textbf{Key Observations}:
\begin{enumerate}
    \item \textbf{Synchronous TX→RX}: The transmission completes on rank 1 before RX notification arrives at rank 2, preserving causality
    \item \textbf{Centralized decisions}: All propagation calculations occur at rank 0, ensuring consistency
    \item \textbf{Transparent to upper layers}: WiFi MAC/PHY layers see standard \texttt{StartReceive()} calls, unaware of MPI
    \item \textbf{Scalable device processing}: While rank 0 calculates propagation, devices on ranks 1-N independently process MAC/PHY protocols
\end{enumerate}

\subsection{Time Synchronization}

Our implementation leverages NS-3's \texttt{DistributedSimulator} for time management. The distributed simulator ensures:

\begin{itemize}
    \item \textbf{Null message protocol}: Ranks periodically exchange messages indicating their local time advancement
    \item \textbf{LBTS calculation}: Each rank computes its Lower Bound on Time Stamp—the earliest possible timestamp of future remote messages
    \item \textbf{Safe event processing}: A rank processes local events only up to its LBTS
    \item \textbf{Lookahead}: WiFi packet transmission times provide natural lookahead (microseconds to milliseconds)
\end{itemize}

This conservative approach guarantees causality without rollback mechanisms, aligning with NS-3's existing MPI framework.

\subsection{Design Decisions and Trade-offs}

\subsubsection{Centralized vs. Distributed Channel}

\paragraph{Decision: Centralized channel processor on rank 0}

\paragraph{Rationale}:
\begin{itemize}
    \item WiFi's broadcast nature resists spatial partitioning
    \item Centralization eliminates consistency problems
    \item Matches NS-3's \texttt{YansWifiChannel} model
    \item Simplifies debugging and validation
\end{itemize}

\paragraph{Trade-off}:
Rank 0 becomes a potential bottleneck. However, propagation calculations are relatively lightweight (microseconds per device pair) compared to device-side processing (full protocol stack). For scenarios with 200 devices and 10 Hz beacons, rank 0 performs ~400K calculations/second—easily handled by modern CPUs.

\subsubsection{Synchronous vs. Asynchronous MPI}

\paragraph{Decision: Synchronous message passing}

\paragraph{Rationale}:
\begin{itemize}
    \item Matches NS-3's discrete event semantics
    \item Simpler implementation (no message buffering complexity)
    \item Causality guaranteed by construction
    \item Aligns with conservative synchronization
\end{itemize}

\paragraph{Trade-off}:
Potential communication latency. Asynchronous MPI could hide latency but requires careful buffering and synchronization. Our measurements (Section~\ref{sec:results}) show synchronous MPI adds acceptable overhead.

\subsubsection{Static vs. Dynamic Device Distribution}

\paragraph{Decision: Static assignment of devices to ranks}

\paragraph{Rationale}:
\begin{itemize}
    \item Simpler implementation (no migration protocol)
    \item Predictable performance characteristics
    \item Sufficient for many scenarios (highway lanes, building floors, network segments)
\end{itemize}

\paragraph{Trade-off}:
Cannot adapt to load imbalance. Future work could explore dynamic redistribution based on runtime profiling.
